\begin{tikzpicture}[
    x=1.25cm,
    y=0.95cm,
    wire/.style={
        draw=diagramfg,
        line cap=round,
        line width=0.8pt
    },
    crossing/.style={
        preaction={
            draw=diagrambg,
            line width=3.5pt,
            line cap=butt
        },
        draw=diagramfg,
        line width=0.8pt,
        line cap=round
    },
    bs/.style={
        draw=none,
        fill=diagramBS,
        rectangle,
        minimum width=5mm,
        minimum height=1.7mm,
        inner sep=0pt,
        line width=0.8pt
    },
    phase/.style={
        draw=diagramfg,
        text=diagramfg,
        fill=diagrambg,
        rectangle,
        rounded corners=1pt,
        minimum width=6mm,
        minimum height=6mm,
        line width=0.8pt,
    },
    meas/.style={
        draw=diagramfg,
        text=diagramfg,
        rectangle,
        rounded corners=1pt,
        minimum width=8mm,
        minimum height=6mm,
        inner sep=0pt,
        line width=0.8pt,
        fill=diagrambg,
        path picture={
            % semicircular dial
            \draw[draw=diagramfg,line width=0.8pt]
                ([xshift=-2.8mm,yshift=-1.6mm]
                    path picture bounding box.center)
                arc[start angle=180,end angle=0,radius=2.8mm];
            % needle with visible arrowhead
            \draw[draw=diagramfg,line width=0.9pt,
                -{Latex[length=1.5mm,width=1.2mm]}
            ]
                ([xshift=0mm,yshift=-1.0mm]
                    path picture bounding box.center)
                --
                ([xshift=1.7mm,yshift=2.6mm]
                    path picture bounding box.center);
        }
    },
    txt/.style={
        font=\small,
        text=diagramfg
    },
    epr/.style={
        draw=diagramaccent,
        dashed,
        line width=0.8pt
    },
    resource label/.style=epr bs/.style={
        draw=diagramfg,
        fill=diagrambg,
        rectangle,
        minimum width=5mm,
        minimum height=5mm,
        inner sep=0pt,
        line width=0.8pt
    },
    classical/.style={
        draw=diagramfg,
        double=diagrambg,
        double distance=1.1pt,
        line width=0.55pt,
        line cap=round,
        line join=round
    },
    classical arrow/.style={
        classical,
        -{Latex[length=1.8mm,width=1.4mm]},
        shorten >=1pt
    },
    disp/.style={
        draw=diagramfg,
        fill=diagrambg,
        circle,
        minimum size=5mm,
        inner sep=0pt,
        line width=0.8pt,
        text=diagramfg,
        font=\small
    }
]

\pgfdeclarelayer{foreground}
\pgfsetlayers{background,main,foreground}

% =================================================
% Rail order: (a,c,c',b,d,d')
% =================================================

\def\ya{5}
\def\yc{4}
\def\ycp{3}
\def\yb{2}
\def\yd{1}
\def\ydp{0}

% =================================================
% Horizontal positions
% =================================================

\def\xLabel{0}
\def\xBoxL{1.45}
\def\xFirstBS{2.25}
\def\xSecondBS{4.0}
\def\xBoxR{4.70}
\def\xPhase{5.50}
\def\xMeasure{6.70}
\def\xResult{7.35}
\def\xOutput{7.45}
\def\xBeforeDisplacement{7.60}
\def\xDisplacement{7.60}
\def\xArrowEnd{8.65}

% short return distances after beamsplitters
\def\xReturnA{2.60}
\def\xReturnB{4.35}



% =================================================
% Left labels
% =================================================

\node[txt,anchor=east] at (\xLabel,\ya) {$\ket{\psi_{\mathrm{in}}}$};
\node[txt,anchor=east] (cLabel) at (\xLabel,\yc) {$\ket{0}_{\hat{x}}$};
\node[txt,anchor=east] (cpLabel) at (\xLabel,\ycp) {$\ket{0}_{\hat{p}}$};
\node[txt,anchor=east] at (\xLabel,\yb) {$\ket{\phi_{\mathrm{in}}}$};
\node[txt,anchor=east] (dLabel) at (\xLabel,\yd) {$\ket{0}_{\hat{x}}$};
\node[txt,anchor=east] (dpLabel) at (\xLabel,\ydp) {$\ket{0}_{\hat{p}}$};

% =================================================
% Straight wires
% =================================================

\draw[wire] (0.12,\ya)  -- (\xBoxL,\ya);

\draw[wire] (0.12,\yb)  -- (\xBoxL,\yb);

% =================================================
% EPR resources
% =================================================
% ---------- upper EPR beamsplitter between c and c' ----------

% left straight parts
\draw[wire] (0.12,\yc)  -- (0.42,\yc);
\draw[wire] (0.12,\ycp) -- (0.42,\ycp);

% gaps
%\draw[crossing] (3.50,\ycp) -- (3.60,\ycp);
%\draw[wire] (0.12,\ycp) -- (0.42,\ycp);


% bend into the BS
\draw[wire] (0.42,\yc)  -- (0.56,3.58);
\draw[wire] (0.42,\ycp) -- (0.56,3.42);

% bend back out to the original rails
\draw[wire] (0.88,3.58) -- (1.02,\yc);
\draw[wire] (0.88,3.42) -- (1.02,\ycp);

% continue on the rails
\draw[wire] (1.02,\yc)  -- (\xBoxL,\yc);
\draw[wire] (1.02,\ycp) -- (\xOutput,\ycp);


% beamsplitter
\node[bs] (EPRup) at (0.72,3.50) {};

% framed label

\node[
    draw=red!75!black,
    rounded corners=2pt,
    line width=0.9pt,
    inner sep=4pt,
    fit=(cLabel)(cpLabel)(EPRup),
    label={[text=red!75!black,font=\small]above:EPR}
    ] (upperEPRbox) {};
% ---------- lower EPR beamsplitter between d and d' ----------

% left straight parts
\draw[wire] (0.12,\yd)  -- (0.42,\yd);
\draw[wire] (0.12,\ydp) -- (0.42,\ydp);


% bend into the BS
\draw[wire] (0.42,\yd)  -- (0.56,0.58) ;
\draw[wire] (0.42,\ydp) -- (0.56,0.42) ;

% continue on the rails
\draw[wire] (1.02,\yd)  -- (\xBoxL,\yd);
\draw[wire] (1.02,\ydp) -- (\xOutput,\ydp);

% bend back out to the original rails
\draw[wire] (0.88,0.58) -- (1.02,\yd);
\draw[wire] (0.88,0.42) -- (1.02,\ydp);

% beamsplitter
\node[bs] (EPRdown) at (0.72,0.50) {};
% framed label
\node[
    draw=red!75!black,
    rounded corners=2pt,
    line width=0.9pt,
    inner sep=4pt,
    fit=(dLabel)(dpLabel)(EPRdown),
    label={[text=red!75!black,font=\small]below:EPR}
] (lowerEPRbox) {};

% =================================================
% First layer: B_ac and B_bd
% =================================================

% Inputs to B_ac
\draw[wire] (\xBoxL,\ya) -- (1.95,\ya) -- (2.12,4.58);
\draw[wire] (\xBoxL,\yc) -- (1.95,\yc) -- (2.12,4.42);

% Inputs to B_bd
\draw[wire] (\xBoxL,\yb) -- (1.95,\yb) -- (2.12,1.58);
\draw[wire] (\xBoxL,\yd) -- (1.95,\yd) -- (2.12,1.42);

% Outputs of B_ac: return to rails a,c
\coordinate (acup) at (\xReturnA,\ya);
\coordinate (aclo) at (\xReturnA,\yc);

\draw[wire] (2.38,4.58) -- (acup);
\draw[wire] (2.38,4.42) -- (aclo);

% Outputs of B_bd: return to rails b,d
\coordinate (bdup) at (\xReturnA,\yb);
\coordinate (bdlo) at (\xReturnA,\yd);

\draw[wire] (2.38,1.58) -- (bdup);
\draw[wire] (2.38,1.42) -- (bdlo);

\node[bs] (Bac) at (\xFirstBS,4.50) {};
\node[bs] (Bbd) at (\xFirstBS,1.50) {};

\node[txt,above=2mm of Bac] {$B_{02}$};
\node[txt,above=2mm of Bbd] {$B_{34}$};

% =================================================
% Second layer: B_ab and B_cd
% =================================================

% Inputs to B_ab: from rails a and b
\draw[wire] (acup) -- (3.15,\ya) -- (3.62,3.88);
\draw[wire] (bdup) -- (3.15,\yb);
\draw[crossing] (3.15,\yb) -- (3.62,3.72);

% Inputs to B_cd: from rails c and d
\draw[wire] (aclo) -- (3.15,\yc);
\draw[crossing] (3.15,\yc) -- (3.62,2.28);
\draw[wire] (bdlo) -- (3.15,\yd) -- (3.62,2.12);

% Outputs of B_ab: return to rails a and b
\coordinate (m1in) at (\xReturnB,\ya);
\coordinate (m3in) at (\xReturnB,\yb);

\draw[wire] (3.92,3.88) -- (m1in);
\draw[crossing] (3.92,3.72) -- (m3in);

% Outputs of B_cd: return to rails c and d
\coordinate (m2in) at (\xReturnB,\yc);
\coordinate (m4in) at (\xReturnB,\yd);

\draw[crossing] (3.92,2.28) -- (m2in);
\draw[wire] (3.92,2.12) -- (m4in);

% continue horizontally to the box edge
\draw[wire] (m1in) -- (\xBoxR,\ya);
\draw[wire] (m2in) -- (\xBoxR,\yc);
\draw[wire] (m3in) -- (\xBoxR,\yb);
\draw[wire] (m4in) -- (\xBoxR,\yd);

% moved a bit higher
\node[bs] (Bab) at (\xSecondBS,3.80) {};
\node[bs] (Bcd) at (\xSecondBS,2.20) {};

\node[txt,above=2mm of Bab] {$B_{03}$};
\node[txt,below=2mm of Bcd] {$B_{14}$};





% =================================================
% Phase rotations
% =================================================

\node[phase] (R3) at (\xPhase,\ya) {$R(\theta_3)$};
\node[phase] (R1) at (\xPhase,\yc) {$R(\theta_1)$};
\node[phase] (R4) at (\xPhase,\yb) {$R(\theta_4)$};
\node[phase] (R2) at (\xPhase,\yd) {$R(\theta_2)$};

\draw[wire] (\xBoxR,\ya) -- (R3.west);
\draw[wire] (\xBoxR,\yc) -- (R1.west);
\draw[wire] (\xBoxR,\yb) -- (R4.west);
\draw[wire] (\xBoxR,\yd) -- (R2.west);

% =================================================
% Measurements
% =================================================

\node[meas] (M3) at (\xMeasure,\ya) {};
\node[meas] (M1) at (\xMeasure,\yc) {};
\node[meas] (M4) at (\xMeasure,\yb) {};
\node[meas] (M2) at (\xMeasure,\yd) {};

\draw[wire] (R3.east) -- (M3.west);
\draw[wire] (R1.east) -- (M1.west);
\draw[wire] (R4.east) -- (M4.west);
\draw[wire] (R2.east) -- (M2.west);


\node[txt,anchor=west] at (\xDisplacement,\ya) {$m_3$};
\node[txt,anchor=west] at (\xDisplacement,\yc) {$m_1$};
\node[txt,anchor=west] at (\xDisplacement,\yb) {$m_4$};
\node[txt,anchor=west] at (\xDisplacement,\yd) {$m_2$};

% =================================================
% Surviving outputs
% =================================================
% Displacement gate on c'
\node[disp] (Dc) at (\xDisplacement,\ycp) {$D$};

% Quantum output wire is split around the gate
\draw[wire] (Dc.east) -- (\xArrowEnd,\ycp);

% Dedicated routing positions
\coordinate (M3exit) at ($(M3.east)+(0.25,0)$);
\coordinate (M1exit) at ($(M1.east)+(0.40,0)$);

% Classical bus above Dc
\coordinate (FFc) at ($(Dc.north)+(0,0.4)$);

% Measurement outcomes
\draw[classical] (M3.east) -- (M3exit) |- (FFc);

\draw[classical] (M1.east) -- (M1exit) |- (FFc);

% Arrow only on the final segment
\draw[classical arrow] (FFc) -- (Dc.north);



% Displacement gate on d'
\node[disp] (Dd) at (\xDisplacement,\ydp) {$D$};

% Quantum output wire
\draw[wire] (Dd.east) -- (\xArrowEnd,\ydp);

% Separate vertical lanes
\coordinate (M4exit) at ($(M4.east)+(0.25,0)$);
\coordinate (M2exit) at ($(M2.east)+(0.40,0)$);

% Classical bus above Dd
\coordinate (FFd) at ($(Dd.north)+(0,0.40)$);

\draw[classical]
    (M4.east) -- (M4exit) |- (FFd);

\draw[classical]
    (M2.east) -- (M2exit) |- (FFd);

\draw[classical arrow]
    (FFd) -- (Dd.north);

\node[
    txt,
    anchor=west
] at ($(\xArrowEnd+0.25,\ycp)$)
    {$\ket{\psi_{\mathrm{out}}}$};

\node[
    txt,
    anchor=west
] at ($(\xArrowEnd+0.25,0)$)
    {$\ket{\phi_{\mathrm{out}}}$};

% =================================================
% Bottom caption
% =================================================

\node[txt,align=center] at (4.85,-0.85)
{
    two-mode gate teleportation:
    $G_{03}(\theta_1,\theta_2,\theta_3,\theta_4)$
};

\end{tikzpicture}
